\documentclass[10pt,a4paper]{amsart}
\usepackage[margin=19mm]{geometry}
\usepackage{amsmath,amssymb,mathtools}
\usepackage[scr=rsfs,cal=euler]{mathalfa}
\usepackage{xfrac}
\usepackage[unicode,hidelinks,bookmarks=false]{hyperref}
\newcommand{\ie}{i.e.\ }
\newcommand{\cf}{cf.\ }
\newcommand{\resp}{respectively\ }
\newcommand{\Subsection}{\subsection}
\providecommand{\nobreakdash}{\nobreak\hbox{-}}
\setlength{\emergencystretch}{2em}
\multlinegap0pt
\usepackage{amssymb}
\usepackage{tikz}
\usepackage{enumerate}
\newtheorem{thm}{Theorem}
\newcommand{\AG}{\text{AG}}
\newcommand{\F}{\mathbb{F}}
\newcommand{\PG}{\text{PG}}
\newcommand{\mpa}[3]{m_{\text{Aff}}(#1 ; #2, #3)}
\title{Large line-free sets and their applications}
\author{Jakob F\"uhrer, Vladislav Taranchuk}
\date{Source: arXiv:2403.18611v2 (CC BY 4.0)}
\begin{document}
\maketitle

\noindent\emph{Source and licence.} Large line-free sets and their applications. Version arXiv:2403.18611v2 (CC BY 4.0), distributed by arXiv under the Creative Commons Attribution 4.0 International licence, http://creativecommons.org/licenses/by/4.0/, as recorded in the arXiv licence metadata for this exact version and shown on the abstract page https://arxiv.org/abs/2403.18611v2. Full licence notice: \texttt{licenses/arxiv-2403.18611-CC-BY-4.0.txt}.\par\smallskip

\noindent\textit{Editorial scope.} Counted unit: the article's thm mainthm (source0.tex lines 149--155). The article's complete original proof of it is delivered with it (source0.tex lines 339--364, one \texttt{proof} environment of 26 lines, which the article places away from the statement). No mathematics is written, repaired or re-derived: the statement and the proof are the article's own lines, in the article's own notation, and no retained line is substituted or re-flowed.  Retained uncounted as the source-local context the proof needs: lines 334--336.  Nothing was omitted from any retained span, so the package carries no omission record.  No retained line contains a citation, so the wrapper carries no bibliography and no bibliography entry.  No retained line contains a figure environment, an image-inclusion command or drawing code, so no image is delivered and the package declares no asset dependency.  Editorial formatting only, in addition: an AMS \texttt{amsart} preamble that restates the article's own declarations and macros used by the retained lines, namely the environments \texttt{thm} on separate counters, together with the standard packages those lines need.  The retained environments print in this document as Theorem 1 in order of appearance, whereas the article prints the same items with the numbering of its own full document.  The complete native source of the article as distributed in the arXiv e-print for this exact version is bundled unchanged as \texttt{sources/156768/source0.tex}.
\begin{equation}\label{NormExpand}
N_d(ax + b) = N_d(a)x^d + tr(ba^{q + q^2 + \cdots +q^{d-1}})x^{d-1} + g_{a, b}(x)    
\end{equation}

\begin{thm}
\label{mainthm}
    Let $q$ be a prime power, and $t < q$. Then 
    $$
    \mpa{t}{n}{q} = \Omega\left(q^{n\left(1  - \frac{2}{t^2 + t}\right)}\right).
    $$
\end{thm}

\begin{proof}[Proof of Theorem \ref{mainthm}]
Let $t \leq q-1$. Identify the vector space $\F_q^{\binom{t+1}{2}}$ with $V = \F_q \times \F_{q^2}\times \cdots \times\F_{q^{t}}$, and so to write $v = (v_1, v_2, \dots, v_{t}) \in V$ means $v_i \in \F_{q^i}$ for each $i$. Now consider the polynomial $P_{t, q}: \F_q^{\binom{t+1}{2}} \mapsto \F_q$ given by
$$
P_{t, q}(v) = N_{t}(v_{t}) + N_{t-1}(v_{t-1}) + \cdots + N_2(v_2) + v_1.
$$
Let $u \in \F_q$, and define 
$$
S_u = \{ v \in V: P_{t, q}(v) = u\}.
$$
\noindent \textbf{Claim}: For any affine line $\ell$ in $V$, $|\ell \cap S_u | \leq t$. 
Suppose for a contradiction that $|\ell \cap S_u| > t$. Parameterize $\ell = \{ ax + b: x\in \F_q \}$ where $a, b \in V$. Then, the polynomial $P_{t, q}(ax + b) - u$
in $x$ has at least $t+1$ roots. Note that by (\ref{NormExpand}), $P_{t, q}(ax + b) - u$ is of degree at most $t$, so if it has more than $t$ roots, then $P_{t, q}(ax + b)$ must be identically 0. Let $a = (a_1, a_2, \dots, a_{t}), b = (b_1, b_2, \dots, b_t)$, and $i \in [t] := \{ 1, 2, \dots, t \}$, be the largest index for which $a_i \neq 0$. This implies that
$$
P_{t, q}(ax + b) - u = N_t(a_tx + b_t) + N_{t-1}(a_{t-1}x + b_{t-1}) + \cdots N_2(a_2x + b_2) + a_1x + b_1-u
$$
has degree $i$ and the coefficient of $x^i$ is precisely $N_i(a_i) \neq 0$. However, if this was the case, it would contradict the fact that $P_{t, q}(ax + b) - u$ is identically 0. Thus $a_i =0$ for each $i$ and so $a = (0, \dots, 0)$, which implies that $\ell = \{ b\}$, a contradiction. Thus, $|\ell \cap S_u| \leq t$ and so $S_u$ is $t$-line evasive and may also be viewed as a $(t,1)$-set in $\PG(\binom{t + 1}{2}, q)$.

Furthermore, we note that $|S_u| = q^{\binom{t+1}{2}- 1}$ since for any choice of $v_2, v_2, \dots, v_{r}$, there exists a unique $v_1$ for which $P_{t, q}((v_1, \dots, v_{t})) = u$.


When $n > \binom{t+1}{2}$, let $k$ be the integer for which the inequality $\binom{t+1}{2}(k-1) <n \leq \binom{t+1}{2}k$ holds. The proof above implies that the set of solutions to $P_{t, q^k}(v) = u$ is $t$-line evasive in $\AG(\binom{t+1}{2}, q^k)$ for each $x \in \F_{q^k}$. Thus, we have a partition of the points of $\AG(\binom{t+1}{2}, q^k)$ into $q^k$ parts such that each part is $t$-line evasive. Performing field reduction yields a partition of $\AG(\binom{t+1}{2}k, q)$ into $q^k$ sets, each of which are $t$-line evasive over $\F_q$. Since $\AG(n, q)$ is a subspace of $\AG(\binom{t+1}{2}k, q)$, then the same holds true for $\AG(n, q)$. Hence, $\AG(n, q)$ contains a $t$-line evasive set of order at least $q^{n-k}$. Since $\binom{t+1}{2}(k-1) < n $, then $k < n/ \binom{t+1}{2} + 1$. Hence, 
$$
\mpa{t}{n}{q} \geq q^{n - k} > q^{n - n/\binom{t+1}{2} - 1} = q^{n\left( 1 - \frac{2}{t^2 + t} \right) - 1}.
$$
When $\binom{t+1}{2}$ divides $n$, then it follows that the set has order precisely $q^{n\left( 1 - \frac{2}{t^2 + t} \right)}$
\end{proof}

\end{document}
